%-----------------------------------------------------------------------
\section{Empirical Results}
\label{sec:results}
%-----------------------------------------------------------------------

\subsection{Latency and Execution Time Analysis}

Figure \ref{fig:keygen_time} outlines absolute execution times across primary key generation bounds. Across non-constrained parameters, classical baselines like ECDH-P256 complete full exchanges significantly under 2ms. During PQC migration, \kyber{}-512 universally dominates across all lattice KEM mechanisms, resolving comprehensively sub-10ms globally. Predictably, \falcon{} executes exceptionally fast verifications yet triggers catastrophic latency barriers internally during generation algorithms owing firmly to complex Fourier bounds unsupported inherently on Cortex-M0 microarchitectures without specialized Floating-Point Units (FPUs). 

\textbf{\sphincs{}} executes as a significant temporal outlier entirely due to thousands of recursive sequential SHA-3 loops. While retaining maximal mathematical conservatism, SLH-DSA key generation times aggressively skew metrics off logarithmic scales within IoT C0 frameworks. 

\begin{figure}[htbp]
\centering
\includegraphics[width=\columnwidth]{fig2_keygen_time}
\caption{Key generation latency plotted logarithmically across simulated architecture groupings. Classical counterparts map consistently closer to zero bounds.}
\label{fig:keygen_time}
\end{figure}

\subsection{Memory Footprint and Scaling Penalty}

Figure \ref{fig:memory} illustrates peak resident memory footprint distribution. A dominant finding isolating "The Class 0 Problem" asserts that Class 0 architecture ($\approx 10$ KB) inherently cannot support \textit{any} existing NIST PQC primitive directly on-chip without overflowing standard program stacks natively triggering hardware fault assertions.

\begin{figure}[htbp]
\centering
\includegraphics[width=\columnwidth]{fig5_memory}
\caption{Peak memory utilization limits measured iteratively against strict hardware allocations.}
\label{fig:memory}
\end{figure}

\subsection{Communication Overhead Analysis}

Unlike classical ECDSA configurations transmitting a meager 64 bytes over RF interfaces, \dilithium{}-3 forces IoT communication busses to manage multiple multi-kilobyte transmissions across fragmented TCP packets.

\begin{itemize}
    \item \textbf{High Payload Cost:} SLH-DSA produces exceptionally large 17 KB ~ 49 KB signatures, triggering massive DTLS packet fragmentations and subsequent datagram dropping. 
    \item \textbf{\falcon{} Efficiency:} Falcon universally outshines competing mechanisms natively when plotting transmission overhead due to generating impressively dense 666-byte signatures. This drastically reduces radio latency albeit demanding heavily on initialization processor stacks.
\end{itemize}

\subsection{Energy Breakdown Assessment}

Correlating Algorithm \ref{alg:energy} calculations yields Figure \ref{fig:energy}. \kyber{}-512 requires significantly less energy ($\approx 4$ mJ on generic C2 bounds) compared to classical mechanisms entirely due to extremely rapid clock-cycle decay minimizing active CPU states. Consequently, LWE KEM mechanisms actually inherently increase embedded battery chronologies when explicitly assuming radio frequency operations (RF tx/rx phases) execute with sufficiently optimized payloads.

\begin{figure}[htbp]
\centering
\includegraphics[width=\columnwidth]{fig12_energy}
\caption{Comprehensive cumulative algorithmic energy expenditures correlated uniformly across target hardware variants.}
\label{fig:energy}
\end{figure}

\subsection{TLS 1.3 Handshake Overheads}

Mapping our derived bounds strictly into the modernized TLS 1.3 sequence geometry \cite{campagna2020tls}, total end-to-end negotiated latency bounds aggressively expand simply processing CertificateVerify sequences padded with full lattice signatures natively transmitted over generic constrained wireless pathways. End-to-end TLS execution bounds realistically skyrocket by factors averaging $\approx 15-25x$ broadly across C1 endpoints exclusively handling fragmented \dilithium{}-2 handshake arrays.
